% BHU PYQ -- Professional Exam Template
\documentclass[11pt,a4paper]{article}

\usepackage[a4paper,top=2.5cm,bottom=2.5cm,left=2.8cm,right=2.8cm,headheight=14pt]{geometry}
\usepackage{amsmath,amssymb}
\usepackage[T1]{fontenc}
\usepackage[utf8]{inputenc}
\usepackage{lmodern}
\usepackage{microtype}
\usepackage{fancyhdr}
\usepackage{enumitem}
\usepackage{hyperref}
\usepackage{lastpage}
\usepackage{parskip}
\usepackage{tikz}
\usetikzlibrary{arrows.meta,shapes,positioning}

\hypersetup{colorlinks=true,urlcolor=black,linkcolor=black,
  pdftitle={B.Sc. (Hons.) Zoology Sem-IV ZOB-401 (Old Course) -- BHU PYQ}}

\pagestyle{fancy}
\fancyhf{}
\renewcommand{\headrulewidth}{0.4pt}
\renewcommand{\footrulewidth}{0.4pt}
\fancyhead[L]{\small\textit{Banaras Hindu University}}
\fancyhead[R]{\small\textit{Zoology Paper: ZOB-401 (Old Course)}}
\fancyfoot[C]{\small Page \thepage\ of \pageref{LastPage}}

\newlist{parts}{enumerate}{3}
\setlist[parts,1]{label=(\a*),leftmargin=*,itemsep=4pt,topsep=3pt}
\setlist[parts,2]{label=(\roman*),leftmargin=*,itemsep=2pt,topsep=2pt}
\setlist[parts,3]{label=(\arabic*),leftmargin=*,itemsep=2pt,topsep=2pt}
\newcommand{\pts}[1]{\hfill\textit{[#1]}}

\begin{document}

\begin{center}
  {\large\bfseries BANARAS HINDU UNIVERSITY}\\[2pt]
  {\normalsize B.Sc. (Hons.) Semester IV Examination, 2013-14 (Old Course)}\\[6pt]
  \rule{0.8\linewidth}{0.5pt}\\[6pt]
  {\large\bfseries Zoology}\\[4pt]
  {\large\bfseries Paper: ZOB-401 --- Genetics, Evolution and Animal Behaviour}\\[4pt]
  {\normalsize Time: Three hours \hfill Full Marks: 70}
\end{center}

\rule{\linewidth}{0.4pt}

\smallskip
\noindent\textit{Note: Answer three questions from Section-A, two questions from Section-B, and two questions from Section-C. Question No. 1 in all sections is compulsory. Use a separate answer book for each section. The figures in the right-hand margin indicate marks.}

\bigskip
\begin{center}
  \textbf{SECTION -- A} \\
  \textbf{(Genetics) (Marks: 36)}
\end{center}

\medskip\noindent\textbf{Question 1.} Answer the following parts:
\begin{parts}
  \item Select the best option in the following:\pts{$2 \times \frac{1}{2} = 1$}
  \begin{parts}
    \item If a homozygous tall plant is crossed with a heterozygous tall plant, which of the following is the most likely outcome for the progeny?
    \begin{parts}
      \item 100 tall : 0 dwarf
      \item 25 tall : 50 dwarf : 25 tall
      \item 75 tall : 25 dwarf
      \item 67 tall : 33 dwarf
    \end{parts}
    
    \item In the ABO blood group system which of the following progeny is not possible?
    \begin{parts}
      \item An A child of A and B parents
      \item An AB child of A and O parents
      \item An O child of A and O parents
      \item An O child of A parents
    \end{parts}
  \end{parts}
  
  \item Write whether the following statements are True or False giving reasons in not more than 2--3 sentences:\pts{$2 \times 3 = 6$}
  \begin{parts}
    \item Drosophila having XY and XXY genotypes are males.
    \item A test cross is performed to find out the genotype of an individual.
    \item Kappa particles are involved in infective inheritance in \textit{Paramecium}.
  \end{parts}

  \item For the terms given in Column-A select the best match from those in Column-B:\pts{$4 \times \frac{1}{2} = 2$}
  \begin{center}
    \begin{tabular}{lp{1cm}l}
      \textbf{Column-A} & & \textbf{Column-B} \\
      (i) Cystic fibrosis & & (1) 45, X \\
      (ii) Down syndrome & & (2) t(9; 22) \\
      (iii) Cry of cat syndrome & & (3) del 5p \\
      (iv) CML & & (4) CFTR \\
    \end{tabular}
  \end{center}

  \item Differentiate between:\pts{$2 \times 1\frac{1}{2} = 3$}
  \begin{parts}
    \item Dominance and epistasis
    \item Penetrance and expressivity.
  \end{parts}
\end{parts}

\medskip\noindent\textbf{Question 2.} Answer the following parts:
\begin{parts}
  \item Define law of segregation and explain its chromosomal basis.\pts{6}
  \item How is G-banding performed and how does this banding help in karyotyping? How are the bands numbered on chromosomes?\pts{6}
\end{parts}

\medskip\noindent\textbf{Question 3.} Answer the following parts:
\begin{parts}
  \item If an individual is heterozygous for a paracentric inversion, what will be the consequences following meiosis?\pts{6}
  \item While carrying out a dihybrid cross the scientist observed a F2 phenotypic ratio of 9:3:4. Explain the phenomenon.\pts{6}
\end{parts}

\medskip\noindent\textbf{Question 4.} Write notes on the following:\pts{$3 \times 4 = 12$}
\begin{parts}
  \item Sex chromosome systems
  \item Genetic counseling
  \item Polygenic inheritance.
\end{parts}

\pagebreak
\begin{center}
  \textbf{SECTION -- B} \\
  \textbf{(Evolution) (Marks: 17)}
\end{center}

\medskip\noindent\textbf{Question 1.} Answer the following parts:
\begin{parts}
  \item Select the best option in the following:\pts{$2 \times \frac{1}{2} = 1$}
  \begin{parts}
    \item During which geological period gigantic dinosaurs, \textit{Archaeopteryx} and early mammals coexisted?
    \begin{parts}
      \item Jurassic
      \item Triassic
      \item Cretaceous
      \item Tertiary
    \end{parts}
    
    \item The earliest ground dwelling hominid was
    \begin{parts}
      \item \textit{Australopithecus}
      \item \textit{Dryopithecus}
      \item \textit{Ramapithecus}
      \item \textit{Homo erectus}
    \end{parts}
  \end{parts}
  
  \item Write whether the following statements are True or False giving reasons in not more than 2--3 sentences:\pts{$2 \times 2 = 4$}
  \begin{parts}
    \item If allelic frequency of A = 60\% and that of a = 40\% in a Mendelian population, then the genotype frequencies will be : AA = 48\%, Aa = 16\% and aa = 36\%.
    \item Higher variation in the sequence of amino acids between different groups of organisms indicates closer phylogenetic relationship.
  \end{parts}
\end{parts}

\medskip\noindent\textbf{Question 2.} Answer the following parts:
\begin{parts}
  \item Explain Natural selection in action, giving a suitable example.\pts{6}
  \item Describe the major trends in human evolution.\pts{6}
\end{parts}

\medskip\noindent\textbf{Question 3.} Write notes on the following:\pts{$3 \times 4 = 12$}
\begin{parts}
  \item Synthetic theory
  \item Significance of homology and analogy in the evolutionary studies
  \item Criticism on Lamarckism.
\end{parts}

\pagebreak
\begin{center}
  \textbf{SECTION -- C} \\
  \textbf{(Animal Behaviour) (Marks: 17)}
\end{center}

\medskip\noindent\textbf{Question 1.} Define the following:\pts{$5 \times 1 = 5$}
\begin{parts}
  \item Innate behaviour
  \item Entrainment
  \item Care-soliciting behaviour
  \item Pheromones
  \item Society.
\end{parts}

\medskip\noindent\textbf{Question 2.} Explain the following:\pts{$2 \times 6 = 12$}
\begin{parts}
  \item Types of learning
  \item Eusocial organization.
\end{parts}

\medskip\noindent\textbf{Question 3.} Write notes on the following:\pts{$4 \times 3 = 12$}
\begin{parts}
  \item Circadian rhythms
  \item Hormones and behaviour
  \item Advantages of group living
  \item Channels of communication.
\end{parts}

\vfill
\rule{\linewidth}{0.4pt}
\begin{center}\small
\href{https://bhu-exam.vercel.app/}{Click Here}\\[4pt]\href{https://me-aryan.vercel.app/}{Click Here}\\[4pt]\href{https://quashan-me.vercel.app/}{Click Here}
\end{center}

\end{document}
